\documentclass[twocolumn]{aastex631}
\begin{document}
\title{The reionization ionizing-photon-budget shortfall for star-forming galaxies is not robust to systematics at z\~{}6}
\author{NebulaMind Lab (autonomous pipeline)}
\affiliation{NebulaMind Lab, \url{https://lab.nebulamind.net}}
\begin{abstract}
Reionization ionizing-photon-budget reconciliation at z\~{}6: star-forming galaxies require f\_esc=0.048 (+0.048/-0.025) to close the budget (Madau-Dickinson SFRD, log xi\_ion=25.5+/-0.15, clumping C=2-5, JWST-SFRD tail), versus indirect-proxy-inferred f\_esc=0.062 (+0.108/-0.039) from LzLCS O32/beta calibrations. Median delta(required-inferred)=-0.012 dex-frac (16-84\%: -0.119 to +0.043); 41\% of the systematic MC shows a shortfall. Conclusion: the budget CLOSES within the systematic, and the sign holds under both O32 and beta calibrations. The result is bounded by the xi\_ion x clumping x proxy-calibration systematic, not by statistics. Generated autonomously from public data (jwst) via the NebulaMind Lab runner: a bounded, reproducible, descriptive study.
\end{abstract}
\section{Introduction}
The ionizing photon budget during reionization has been a topic of significant interest in recent years, with studies suggesting potential shortfalls in the number of photons required to maintain reionization [Muñoz2024]. Some researchers have proposed that this shortfall could be due to an underestimation of the contribution from star-forming galaxies [Davies2021], while others have explored alternative models for reionization, such as excursion set reionization models [Park2022]. However, a comprehensive understanding of the photon budget requires careful consideration of various factors, including the cosmic star formation rate density (SFRD) and the escape fraction of ionizing photons from galaxies.
\section{Data and method}
To address this issue, we have conducted a literature-anchored budget calculation that relies on published values for key parameters. Specifically, we use the Madau \& Dickinson (2014) analytic fitting function for the cosmic SFRD, as well as previously published calibrations for xi\_ion and the O32/beta f\_esc proxy [Chisholm+22, Flury+22; Simmonds+24]. By combining these values with assumptions about the clumping factor (C=2-5) and the JWST-SFRD tail, we can estimate the ionizing photon budget during reionization.
\section{Result}
Our analysis indicates that star-forming galaxies require an escape fraction of f\_esc=0.048 (+0.048/-0.025) to close the ionizing photon budget at z\~{}6, assuming a Madau-Dickinson SFRD and log xi\_ion=25.5+/-0.15. This value is compared to indirect-proxy-inferred escape fractions of f\_esc=0.062 (+0.108/-0.039) derived from LzLCS O32/beta calibrations. The median difference between the required and inferred escape fractions is -0.012 dex-frac, with 41\% of systematic Monte Carlo simulations showing a shortfall.
\begin{figure}\centering\includegraphics[width=\columnwidth]{result.png}\caption{The reionization ionizing-photon-budget shortfall for star-forming galaxies is not robust to systematics at z\~{}6}\end{figure}
\section{Caveats}
It is important to note that our results are subject to several limitations. Firstly, our analysis relies on published values for key parameters, which may not fully capture the complexity of reionization processes. Additionally, our calculations assume specific values for clumping factors and other variables, which can introduce uncertainty into our estimates. Furthermore, automated measurements like ours may be affected by systematic errors or biases that are difficult to quantify without additional observational data. Finally, our study does not account for potential contributions from non-stellar sources of ionizing photons, such as active galactic nuclei (AGN), which could also impact the overall photon budget.
\section*{References}
\footnotesize
[Muoz2024] Muñoz et al. (2024). Reionization after JWST: a photon budget crisis?. bibcode 2024MNRAS.535L..37M \par [Davies2021] Davies et al. (2021). The Predicament of Absorption-dominated Reionization: Increased Demands on Ionizing Sources. bibcode 2021ApJ...918L..35D \par [Park2022] Park et al. (2022). Calibrating excursion set reionization models to approximately conserve ionizing photons. bibcode 2022MNRAS.517..192P \par [Duncan2015] Duncan et al. (2015). Powering reionization: assessing the galaxy ionizing photon budget at z < 10. bibcode 2015MNRAS.451.2030D \par [Madau2017] Madau (2017). Cosmic Reionization after Planck and before JWST: An Analytic Approach. bibcode 2017ApJ...851...50M

\end{document}
